\documentclass[11pt]{amsart}

\usepackage{amsmath,amssymb,amsthm,mathtools}
\usepackage[a4paper,margin=27mm]{geometry}
\usepackage[hidelinks]{hyperref}

\newtheorem{theorem}{Theorem}[section]
\newtheorem{proposition}[theorem]{Proposition}
\newtheorem{lemma}[theorem]{Lemma}
\theoremstyle{remark}
\newtheorem{remark}[theorem]{Remark}

\newcommand{\Z}{\mathbb Z}
\newcommand{\Q}{\mathbb Q}
\newcommand{\Sym}{\operatorname{Sym}}
\newcommand{\Hom}{\operatorname{Hom}}
\newcommand{\gm}{\gamma}
\newcommand{\aug}{\varpi}
\newcommand{\Rel}{\mathcal R}
\newcommand{\pr}{\operatorname{pr}}

\title[Odd-dimensional vanishing for finite metabelian Lie rings]
{A PBW and derived-functor proof of an odd-dimensional vanishing theorem\\
for finite metabelian Lie rings}
\author{}
\date{}

\begin{document}

\begin{abstract}
Let $L$ be a finite metabelian Lie ring of nilpotency class at most
$n+1$ and suppose that its last nonzero lower-central term
$\gamma_{n+1}(L)$ is a cyclic $2$-group.  We prove that
$\delta_{2n+1}(L)=0$.  The proof uses PBW collection and the ordinary
Koszul complexes for derived symmetric powers.  The intermediate PBW
packets need not be cycles; instead, a chain-level transgression kills
their contributions directly.  The last, quadratic, packet is a genuine
cycle and is annihilated by an integral Bockstein certificate.  At no
point is a homogeneous tail of a relation treated as a relation.
\end{abstract}

\maketitle

\section{Statement and preliminary reductions}

For a Lie ring $L$, let $U(L)$ be its universal enveloping algebra and
$\aug(L)$ its augmentation ideal.  Put
\[
 \delta_r(L)=L\cap\aug(L)^r,
 \qquad
 \gm_1(L)=L,
 \qquad
 \gm_{r+1}(L)=[\gm_r(L),L].
\]
All Lie rings in this paper are Lie rings over $\Z$.

\begin{theorem}\label{thm:main}
Let $n\geq1$, and let $L$ be a finite metabelian Lie ring such that
$\gm_{n+2}(L)=0$ and
\[
 C:=\gm_{n+1}(L)\cong\Z/2^e
\]
is cyclic.  Then
\[
                         \delta_{2n+1}(L)=0.
\]
\end{theorem}

Whenever $C=\langle z\rangle\cong\Z/q$ occurs below, we identify it
with $\frac1q\Z/\Z\subseteq\Q/\Z$ by $z\mapsto\frac1q+\Z$.

We use the following two standard facts.  For every Lie ring,
$\delta_3(L)=\gm_3(L)$.  For every metabelian Lie ring,
\begin{equation}\label{eq:PS}
                         2\delta_r(L)\subseteq\gm_r(L).
\end{equation}
The latter assertion is due to Passi and Sicking
\cite[Theorem~2.1]{PS}.

\begin{lemma}[cyclic last-layer reduction]\label{lem:cyclic-reduction}
Fix $s\geq2$.  Suppose Theorem~\ref{thm:main} is known with $s$ in
place of $n$ whenever $\gm_{s+1}$ is a cyclic $2$-group, and suppose
the theorem is already known in every class at most $s$.  Then it is
valid, with index $2s+1$, for every finite metabelian Lie ring of
class at most $s+1$.
\end{lemma}

\begin{proof}
Let $A$ be such a Lie ring and let $0\ne a\in\delta_{2s+1}(A)$.
Since $\gm_{2s+1}(A)=0$, \eqref{eq:PS} gives $2a=0$.  The image of
$a$ in $A/\gm_{s+1}(A)$ belongs to
$\delta_{2s+1}\subseteq\delta_{2s-1}=0$ by the hypothesis in lower
class.  Hence $a$ lies in the central group $\gm_{s+1}(A)$.  Its
$2$-primary part is a finite abelian $2$-group.
There is a homomorphism from that group to a cyclic $2$-group which is
nonzero on $a$: write the group as a direct sum of cyclic groups and
project onto a summand on which $a$ has nonzero coordinate.  Extend
the kernel by the odd-primary part.  This kernel is an ideal because
$\gm_{s+1}(A)$ is central.  In the resulting quotient the image of
$a$ is nonzero, still belongs to $\delta_{2s+1}$, and the last
lower-central term is cyclic of $2$-power order, contradicting the
hypothesis.
\end{proof}

\begin{proposition}[reduction to the top layer]\label{prop:top-reduction}
To prove Theorem~\ref{thm:main} it is enough to show that, for
$C=\langle z\rangle\cong\Z/q$, $q=2^e$,
\begin{equation}\label{eq:top-reduction}
 \zeta z\in\delta_{2n+1}(L)
 \quad\Longrightarrow\quad
 \zeta\equiv0\pmod q.
\end{equation}
\end{proposition}

\begin{proof}
We argue by induction on $n$.  If $n=1$, then
$\delta_3(L)=\gm_3(L)=0$.  Assume the assertion known for $n-1$.
For $n=2$, the equality $\delta_3(A)=\gm_3(A)$ gives
$\delta_3(A)=0$ for every $A$ of class at most two.  For $n>2$,
Lemma~\ref{lem:cyclic-reduction} gives
$\delta_{2n-1}(A)=0$ for every finite metabelian Lie ring $A$ of
class at most $n$.  If $a\in\delta_{2n+1}(L)$, then its image in
$L/C$ belongs to
\[
 \delta_{2n+1}(L/C)\subseteq\delta_{2n-1}(L/C)=0.
\]
Thus $a\in C$ and has the form $a=\zeta z$.  Now
\eqref{eq:top-reduction} gives $a=0$.
\end{proof}

\section{The graded free model and the PBW witness}

Choose a finite generating set of $L$ and write
\[
                         L=F/\Rel,
\]
where $F$ is the relatively free metabelian Lie ring of class
$n+1$ on a finitely generated free abelian group $X$.  The usual Hall
basis gives a grading by bracket weight
\[
 F=\bigoplus_{s=1}^{n+1}F_s,
 \qquad F_1=X,
 \qquad [F_i,F_j]\subseteq F_{i+j}.
\]
Every $F_s$ is free abelian.  Put $M=F_{\geq2}=[F,F]$.  Metabelianness
gives
\begin{equation}\label{eq:metabelian-grading}
 [F_i,F_j]=0\quad(i,j\geq2),
 \qquad
 [F_i,F_1]\subseteq F_{i+1}.
\end{equation}
In particular $M$ is a module over $\Sym(X)$ by
$m\cdot x=[m,x]$; Jacobi and $[M,M]=0$ make this action symmetric.

Let $U=U(F)$.  Order a homogeneous basis of $F$ first by weight and
then arbitrarily inside each weight.  PBW identifies $U$, as an
abelian group, with the free group on the corresponding ordered
monomials.  Give a PBW monomial total weight equal to the sum of the
weights of its Lie factors, and denote the weight-$w$ subgroup by
$U_w$.

\begin{lemma}\label{lem:augmentation-weight}
For every $r\geq1$,
\[
                         \aug(F)^r=\bigoplus_{w\geq r}U_w.
\]
\end{lemma}

\begin{proof}
Every element of $F_s$ is a bracket of length $s$ in $X$, hence lies
in $\aug(F)^s$.  Therefore every PBW monomial of total weight $w$
belongs to $\aug(F)^w$.  This proves the inclusion from right to left.
Conversely, $\aug(F)$ is generated by homogeneous elements of positive
weight, so a product of $r$ of its elements has weight at least $r$.
\end{proof}

Let $\mathfrak r$ be the two-sided ideal of $U$ generated by $\Rel$.
Since
\[
 x\rho=\rho x+[x,\rho]
 \qquad(x\in F,\ \rho\in\Rel)
\]
and $[x,\rho]\in\Rel$, one has
\begin{equation}\label{eq:left-ideal}
                         \mathfrak r=\Rel U.
\end{equation}
Fix $\widetilde z\in F_{n+1}$ mapping to $z$.  If
$\zeta z\in\delta_{2n+1}(L)$, then there is
\begin{equation}\label{eq:witness}
 \theta\in\Rel U,
 \qquad
 \theta-\zeta\widetilde z\in\aug(F)^{2n+1}.
\end{equation}
Write $\theta_w$ for the weight-$w$ part.
Lemma~\ref{lem:augmentation-weight} gives
\begin{equation}\label{eq:governing}
 \theta_w=0\quad(w\leq2n,\ w\ne n+1),
 \qquad
 \theta_{n+1}=\zeta\widetilde z.
\end{equation}
Moreover, within each fixed weight the PBW summands with different
numbers of Lie factors are independent.  Thus every multi-factor
PBW coefficient of $\theta$ of weight at most $2n$ is zero.
Since a one-factor PBW monomial has weight at most $n+1$, one also has
\begin{equation}\label{eq:primitive-witness}
                         \pr_1(\theta)=\zeta\widetilde z,
\end{equation}
where $\pr_1:U\to F$ denotes the one-factor PBW projection.

\section{Canonical presentations and Koszul complexes}

For $2\leq k\leq n$, put
\[
 U_k=\gm_k(L)/\gm_{k+1}(L),\qquad
 W_k=L/\gm_{k+1}(L),\qquad
 V_k=L/\gm_k(L).
\]
There is an exact sequence of additive groups
\begin{equation}\label{eq:tower-extension}
 E_k:\qquad 0\longrightarrow U_k\longrightarrow W_k
 \xrightarrow{\pi_k}V_k\longrightarrow0.
\end{equation}

There is a convenient presentation obtained directly from $F$.  Set
\[
 A_k=F_{\leq k},
 \qquad
 D_k=\pr_{\leq k}(\Rel)\subseteq A_k.
\]
Then
\begin{equation}\label{eq:canonical-resolution}
 0\longrightarrow D_k\xrightarrow{d_k}A_k
 \longrightarrow W_k\longrightarrow0
\end{equation}
is exact, where $d_k$ is inclusion.  Projection onto weights at most
$k-1$ gives a morphism from \eqref{eq:canonical-resolution} to the
analogous resolution
\[
 0\longrightarrow D_{k-1}\xrightarrow{d_{k-1}}A_{k-1}
 \longrightarrow V_k\longrightarrow0.
\]
Indeed, the kernel of $A_k\to L/\gm_{k+1}(L)$ is exactly
$\pr_{\leq k}(\Rel)=D_k$.

For a two-term free resolution $P=(P_1\xrightarrow dP_0)$, let
\begin{equation}\label{eq:koszul}
 K_m(P)_i=\Lambda^iP_1\otimes\Sym^{m-i}P_0
 \qquad(0\leq i\leq m)
\end{equation}
with the usual Koszul differential
\[
 \partial(a_1\wedge\cdots\wedge a_i\otimes u)
 =\sum_{j=1}^i(-1)^{i-j}
 a_1\wedge\cdots\widehat{a_j}\cdots\wedge a_i
 \otimes da_j\,u.
\]
Then $H_iK_m(P)=L_i\Sym^m(\operatorname{coker}d)$; see, for example,
\cite{BM}.

We use only these full complexes.  Write
\[
 K_m(k)=K_m(D_k\longrightarrow A_k).
\]
Projection onto weights at most $k-1$ induces a chain map
\[
 f_k:K_m(k)\longrightarrow K_m(k-1).
\]

\section{The higher characters and their chain contractions}

Put
\[
                         m_k=n+2-k.
\]
The bracket gives a well-defined symmetric map
\begin{equation}\label{eq:theta-k}
 \Theta_k:U_k\otimes\Sym^{m_k-1}(V_k)\longrightarrow C,
 \qquad
 \Theta_k(\bar u;x_1\cdots x_{m_k-1})
 =[u,x_1,\ldots,x_{m_k-1}].
\end{equation}
Indeed, changing $u$ by an element of $\gm_{k+1}$ produces a bracket
in $\gm_{n+2}=0$; changing one $x_i$ by an element of $\gm_k$ gives
a bracket of two elements of $L'$.  Symmetry follows from Jacobi and
$[L',L']=0$.

The extension cocycle of \eqref{eq:tower-extension} can be described
on the canonical resolutions.  If $a\in D_{k-1}$, choose
$\rho\in\Rel$ with $\pr_{\leq k-1}\rho=a$ and put
\[
 b_k(a)=-\rho_k\pmod{D_k\cap F_k}\in U_k.
\]
The result is independent of the choice of $\rho$.  On the
degree-one part of $K_{m_k}(D_{k-1}\to A_{k-1})$, define
\begin{equation}\label{eq:phi-k}
 \Phi_k(a;x_1\cdots x_{m_k-1})
 =-\Theta_k(b_k(a);\bar x_1\cdots\bar x_{m_k-1}).
\end{equation}
It is a cocycle: in the value on a Koszul boundary one argument is
$d_{k-1}a'$, whose image in $V_k$ is zero.  It is supported in total
bracket weight $n+1$.  Consequently it defines the natural character
\begin{equation}\label{eq:higher-eta}
 \eta_k:L_1\Sym^{m_k}(V_k)\longrightarrow\Q/\Z,
 \qquad \eta_k([c])=\Phi_k(c)
 \quad(2\leq k<n).
\end{equation}
This character depends only on the extension $E_k$ and $\Theta_k$.
Indeed, replacing the cocycle lift $b_k$ by
$b_k+\lambda d_{k-1}$ changes $\Phi_k$ by the coboundary of
\[
 S_k(x_1\cdots x_{m_k})
 =-\sum_i\Theta_k(\lambda x_i;
       \bar x_1\cdots\widehat{\bar x_i}\cdots\bar x_{m_k}).
\]
The usual comparison maps between free resolutions then show that
$\eta_k$ is natural.

For $k<n$ there is a chain-level null-homotopy after pullback.  Write
$w=w_Q+w_P$ according to $A_k=A_{k-1}\oplus F_k$, and denote the
images of these components in $V_k$ and $U_k$ by $\bar w_Q$ and
$\bar w_P$.  Define
\begin{equation}\label{eq:T-k}
 T_k(w_1\cdots w_{m_k})=
 \sum_{i=1}^{m_k}
 \Theta_k(\bar w_{i,P};
     \bar w_{1,Q}\cdots\widehat{\bar w_{i,Q}}\cdots
     \bar w_{m_k,Q}).
\end{equation}

\begin{lemma}[transgression]\label{lem:transgression}
For $2\leq k<n$,
\[
                         f_k^*\Phi_k=\delta T_k.
\]
In particular,
\begin{equation}\label{eq:chain-transgression}
 \Phi_k(f_kc)=T_k(\partial c)
 \qquad(c\in K_{m_k}(k)_1),
\end{equation}
and
\begin{equation}\label{eq:higher-vanishing}
 \eta_k\circ L_1\Sym^{m_k}(\pi_k)=0.
\end{equation}
\end{lemma}

\begin{proof}
Let $r\in D_k$ and $w_2,\ldots,w_{m_k}\in A_k$.  The
$F_k$-component of $d_kr$ is the negative of the cocycle tail
$b_k(\pr_{\leq k-1}r)$, modulo a relation in $U_k$.  In
$T_k((d_kr)w_2\cdots w_{m_k})$, the summand in which the first factor
occupies the $U_k$-slot is therefore
$f_k^*\Phi_k(r;w_2\cdots w_{m_k})$.  Every other summand
contains the image of $d_kr$ in $V_k$ and is zero.  This proves the
identity and hence \eqref{eq:chain-transgression}.  Evaluation on a
cycle gives \eqref{eq:higher-vanishing}.  Notice also that $T_k$ can
be nonzero only when its $F_k$-argument has weight $k$ and all other
arguments have weight one: if another argument has weight at least
two, \eqref{eq:metabelian-grading} makes the long bracket zero.  Thus
$T_k$ sees precisely total bracket weight
$k+(m_k-1)=n+1$.  This support statement is all that the PBW argument
will need.
\end{proof}

\section{The quadratic character}

We now treat $k=n$, so $m_n=2$.  Write
\[
 0\longrightarrow U\longrightarrow W\xrightarrow\pi V
 \longrightarrow0
\]
for $E_n$, and retain $C=\Z/q$.  Choose Smith free resolutions
$P_1\xrightarrow dP_0\to V$ and $Q_1\xrightarrow eQ_0\to U$, and a
lift $\widetilde b:P_1\to Q_0$ of the extension cocycle.  A resolution
of $W$ is
\begin{equation}\label{eq:quadratic-resolution}
 R_1=P_1\oplus Q_1\xrightarrow\partial R_0=P_0\oplus Q_0,
 \qquad
 \partial(a,s)=(da,es-\widetilde b(a)).
\end{equation}

Let $g:V\otimes U\to C$ be $g(x,u)=[u,x]$.  Define
\begin{equation}\label{eq:quadratic-phi}
 \varphi(a,x)=g(\bar x,\overline{\widetilde b(a)}).
\end{equation}
This is a cocycle on $K_2(P)$, since the image of $da'$ in $V$ is zero.
Choose an integral lift $\widehat\varphi$.  It may be chosen so that
\begin{equation}\label{eq:exterior-lift}
 \widehat\varphi(a,da')+\widehat\varphi(a',da)=0.
\end{equation}
Indeed, write $da_i=d_ix_i$, with $d_i\mid d_j$ for $i<j$, and lift
$x_i$ to $X_i\in L$.  The extension cocycle is represented by
$d_iX_i\in U$, so
\[
 \varphi(a_i,x_j)=d_i[X_i,X_j],
 \qquad
 \varphi(a_j,x_i)=-d_j[X_i,X_j]
 =-\frac{d_j}{d_i}\varphi(a_i,x_j)
\]
in $C$.  Choose $r_{ij}\in\Z$ with
$r_{ij}/q+\Z=\varphi(a_i,x_j)$, and set
\[
 \widehat\varphi(a_i,x_j)=r_{ij},
 \qquad
 \widehat\varphi(a_j,x_i)=-\frac{d_j}{d_i}r_{ij},
 \qquad
 \widehat\varphi(a_i,x_i)=0.
\]
This gives
\eqref{eq:exterior-lift} exactly over $\Z$.  The Koszul homology and the
resulting character are independent of the chosen free resolution.

Because $V$ is finite, $d\otimes\Q$ is an isomorphism and
$H_2K_2(P)=0$.  Universal coefficients give
\begin{equation}\label{eq:UCT}
 H^2\Hom(K_2(P),\Z)
 \cong\Hom(L_1\Sym^2(V),\Q/\Z).
\end{equation}
The integral cochain
\[
 h_{\widehat\varphi}(a\wedge a')
 =\frac{\widehat\varphi(a,da')}{q}
\]
is well defined because $da'$ represents zero in $V$, so the numerator
is divisible by $q$.  It
defines, by \eqref{eq:UCT}, a character
\begin{equation}\label{eq:eta}
                         \eta_n:L_1\Sym^2(V)\to\Q/\Z.
\end{equation}
If $\beta_S$ is the Bockstein associated with
$0\to\Z\xrightarrow q\Z\to C\to0$, then
\begin{equation}\label{eq:half}
                         \beta_S[\varphi]=2\eta_n.
\end{equation}
Indeed, on $a\wedge a'$ its integral representative is
\[
 \frac{\widehat\varphi(a,da')-\widehat\varphi(a',da)}q
 =2h_{\widehat\varphi}(a\wedge a')
\]
by \eqref{eq:exterior-lift}.
Thus $\eta_n$, rather than its double
$\beta_S[\varphi]=2\eta_n$, is the quadratic PBW character.

\begin{proposition}[canonical quadratic PBW block]
\label{prop:quadratic-block}
Suppose that \eqref{eq:quadratic-resolution} is realized by the graded
presentation $F\to L$ used above.  Choose Smith bases
\[
 da_i=d_ix_i,
 \qquad es_\alpha=e_\alpha y_\alpha,
 \qquad d_i\mid d_j\quad(i<j),
\]
put $B_i=\widetilde b(a_i)$, and order every $x_i$ before every
$y_\alpha$ in the PBW basis.  Choose full relations
\begin{equation}\label{eq:quadratic-full-relations}
 \rho_i=d_ix_i-B_i+t_i,
 \qquad
 \sigma_\alpha=e_\alpha y_\alpha+t'_\alpha,
 \qquad t_i,t'_\alpha\in F_{n+1}.
\end{equation}
Let
\[
 \chi=c+v+v'+w\in
 (P_1\otimes P_0)\oplus(P_1\otimes Q_0)
 \oplus(Q_1\otimes P_0)\oplus(Q_1\otimes Q_0)
\]
be a cycle.  It has a canonical placed realization obtained as follows:
write
\begin{align*}
 c&=\sum_{i<j}u_{ij}
 \left(\frac{d_j}{d_i}a_i\otimes x_j-a_j\otimes x_i\right),\\
 v&=\sum v_{i\alpha}a_i\otimes y_\alpha,
 &v'&=\sum v'_{\alpha i}s_\alpha\otimes x_i,
 &w&=\sum w_{\alpha\beta}s_\alpha\otimes y_\beta,
\end{align*}
and use the placed rows
\begin{equation}\label{eq:canonical-quadratic-rows}
 \begin{split}
 &u_{ij}\left(\frac{d_j}{d_i}\rho_i x_j-x_i\rho_j\right),
 \qquad v_{i\alpha}\rho_i y_\alpha,\\
 &v'_{\alpha i}x_i\sigma_\alpha,
 \qquad w_{\alpha\beta}\sigma_\alpha y_\beta.
 \end{split}
\end{equation}
If the remaining one-factor coefficient of this complete block in $C$
is $\kappa z$, then
\begin{equation}\label{eq:quadratic-PBW-block}
 \frac\kappa q+\Z
 =-\eta_n\bigl(L_1\Sym^2(\pi)[\chi]\bigr).
\end{equation}
Every placed two-factor block with vanishing two-factor coefficients is
carried to \eqref{eq:canonical-quadratic-rows} by integral relation
interchanges; their correction terms are full one-factor relations.
\end{proposition}

\begin{proof}
First, a cycle in the $P_1\otimes P_0$ summand has the displayed form.
Indeed, the coefficient of $x_i^2$ is $d_i$ times the diagonal
coefficient and hence forces that coefficient to be zero; the
$x_ix_j$ coefficient gives
$d_i c_{ij}+d_jc_{ji}=0$, and $d_i\mid d_j$.  The remaining components
of $\partial\chi=0$, sorted in
\[
 \Sym^2(P_0\oplus Q_0)=
 \Sym^2P_0\oplus(P_0\otimes Q_0)\oplus\Sym^2Q_0,
\]
are
\begin{align}
 (d\otimes1)c&=0,\label{eq:block-A}\\
 -(\widetilde b\otimes1)c
 +(d\otimes1)v+(e\otimes1)v'&=0,\label{eq:block-B}\\
 -\mu(\widetilde b\otimes1)v+(e\otimes1)w&=0.
 \label{eq:block-C}
\end{align}
Thus the two-factor PBW symbol of \eqref{eq:canonical-quadratic-rows}
is zero.  Conversely, an adjacent interchange
$x\rho=\rho x+[x,\rho]$ changes no symmetric two-factor symbol and
adds a full one-factor relation.  It therefore converts any placement
with this symbol, integrally, to the canonical one.

We now calculate its primitive term, rather than infer it from
\eqref{eq:block-A}--\eqref{eq:block-C}.  In the horizontal $i,j$ row,
the $P_0P_0$ products cancel before any interchange:
\[
 \frac{d_j}{d_i}d_ix_ix_j-x_id_jx_j=0.
\]
The $F_n$-tail of its first relation gives
\[
 -\frac{d_j}{d_i}B_ix_j
 =-\frac{d_j}{d_i}x_jB_i
  -\frac{d_j}{d_i}[B_i,x_j],
\]
whereas $x_iB_j$ in the second placement is already ordered.  Hence
the horizontal primitive coefficient is
\begin{equation}\label{eq:block-Z}
 \kappa\equiv-
 \sum_{i<j}u_{ij}\frac{d_j}{d_i}
 \widehat\varphi(a_i,x_j)\pmod q.
\end{equation}
There are no suppressed vertical corrections.  In $\rho_i y_\alpha$
the only possibly inverted tail product is $B_i y_\alpha$, and its
commutator is zero because both factors lie in the metabelian derived
ideal.  The products $x_i\sigma_\alpha$ are already ordered, and the
$\sigma_\alpha y_\beta$ products again have two derived factors.
Finally, the $t_i,t'_\alpha$ in
\eqref{eq:quadratic-full-relations} have weight $n+1$, so multiplying
them by a positive-weight factor cannot produce a primitive term in
weight $n+1$.  This proves \eqref{eq:block-Z} coefficient by
coefficient.

On the other hand,
\[
 \partial(a_i\wedge a_j)=d_i
 \left(\frac{d_j}{d_i}a_i\otimes x_j-a_j\otimes x_i\right),
 \qquad
 h_{\widehat\varphi}(a_i\wedge a_j)
 =\frac{d_j\widehat\varphi(a_i,x_j)}q.
\]
The universal-coefficient description \eqref{eq:UCT} gives
\[
 \eta_n([c])=
 \frac1q\sum_{i<j}u_{ij}\frac{d_j}{d_i}
 \widehat\varphi(a_i,x_j)+\Z.
\]
Projection of $[\chi]$ along $\pi$ is $[c]$; comparison with
\eqref{eq:block-Z} proves \eqref{eq:quadratic-PBW-block}.
\end{proof}

\subsection*{The certificate}

Put $A=\Hom(U,C)$ and $A^\vee=\Hom(A,C)\cong U/qU$, where the last
isomorphism is the evaluation map (it is an isomorphism on each cyclic
summand of $U$).  Let
\[
 \rho:P_0\to A,
 \quad \rho(x)=g(\bar x,-),
 \qquad
 b:P_1\to A^\vee,
 \quad b(a)=\overline{\widetilde b(a)}.
\]
Form the pushout
\[
 \mathcal E=(A^\vee\oplus P_0)/
 \langle(b(a),-da):a\in P_1\rangle.
\]
The map $\mathcal E\to V$, $[\chi,x]\mapsto\bar x$, is onto and its
kernel is generated by the finite group $A^\vee$; hence $\mathcal E$ is
finite.
Extend $\widehat\varphi$ rationally and define the alternating form
$\vartheta$ on $P_0\otimes\Q$ by
\[
 \vartheta(da,x)=\widehat\varphi(a,x)/q.
\]
Equation \eqref{eq:exterior-lift} makes it alternating.  On
$\mathcal E$ define
\[
 \Omega([\chi,x],[\chi',x'])
 =\chi(\rho x')-\chi'(\rho x)+\vartheta(x,x')\pmod\Z.
\]
It is well defined: substituting $(b(a),-da)$ in either variable
gives
$\varphi(a,x)-\widehat\varphi(a,x)/q=0$ in $\Q/\Z$.

Every alternating form on a finite abelian group with values in the
divisible group $\Q/\Z$ is the antisymmetrization of a bilinear form:
choose cyclic generators and place its matrix in one triangle.  Choose
such a $\Q/\Z$-valued form $H$ with
$H^{\mathsf T}-H=\Omega$.  Although $H$ need not be $C$-valued on all
of $\mathcal E\times\mathcal E$, its value is $C$-valued whenever one
argument lies in $i(A^\vee)$: every element of $A^\vee\cong U/qU$ is
killed by $q$, and bilinearity then kills the corresponding value by
$q$.
If
$i(\chi)=[\chi,0]$ and $j(x)=[0,x]$, define
\[
 \chi(\alpha x)=H(jx,i\chi),
 \qquad
 \lambda(\chi,\chi')=H(i\chi',i\chi).
\]
The preceding observation and the evaluation duality
$\Hom(A^\vee,C)\cong A$ show that $\alpha:P_0\to A$ is well defined.
Also $\lambda$ is $C$-valued and symmetric, and the relation
$jd=ib$ gives
\begin{equation}\label{eq:certificate-relation}
                         \alpha d=\lambda_*b.
\end{equation}
Indeed, for $\chi\in A^\vee$,
\[
 \chi(\alpha da)=H(jda,i\chi)=H(ib(a),i\chi)
 =\lambda(\chi,b(a)).
\]
Moreover, the cocycle
\[
                         \psi(a,x)=b(a)(\alpha x)
\]
satisfies
\begin{equation}\label{eq:certificate-bockstein}
                         \beta_S[\psi]=\eta_n.
\end{equation}
It is a cocycle because
\[
 \psi(a,da')-\psi(a',da)
 =\lambda(b(a),b(a'))-\lambda(b(a'),b(a))=0.
\]
To compute its Bockstein, lift
$(x,x')\mapsto H(jx,jx')$ to a rational bilinear form
$P$ on the free group $P_0$.  The restriction
$H(jP_0,jdP_1)=H(jP_0,ib(P_1))$ is $C$-valued, so every such lift
satisfies $qP(x,da)\in\Z$.  Adding an integral bilinear form in one
triangle, arrange $P^{\mathsf T}-P=\vartheta$ without changing that
integrality.  Then $\widehat\psi(a,x)=qP(x,da)$ is integral, reduces to
$\psi$, and
\[
 \widehat\psi(a,da')-\widehat\psi(a',da)
 =\widehat\varphi(a,da').
\]
After division by $q$, this is exactly
\eqref{eq:certificate-bockstein}.

\begin{proposition}[quadratic vanishing]\label{prop:quadratic-vanishing}
\[
                         \eta_n\circ L_1\Sym^2(\pi)=0.
\]
\end{proposition}

\begin{proof}
On $R_0=P_0\oplus Q_0$ define the symmetric $C$-valued form
\[
 \begin{split}
 T((x,y),(x',y'))={}&-\alpha(x)(y'_U)
                     -\alpha(x')(y_U)\\
                    &-\lambda(\bar y,\bar y'),
 \end{split}
\]
where $y_U$ denotes the image of $y$ in $U$ and $\bar y$ its image in
$U/qU$.  Using
\eqref{eq:quadratic-resolution} and
\eqref{eq:certificate-relation}, one obtains
\[
 \begin{aligned}
 T(\partial(a,s),(x,y))
 &=-\lambda(b(a),\bar y)+b(a)(\alpha x)
   +\lambda(b(a),\bar y)\\
 &=\psi(a,x).
 \end{aligned}
\]
Thus the pullback of $[\psi]$ to $K_2(R)$ is zero.  Naturality of the
Bockstein and \eqref{eq:certificate-bockstein} give
\[
 \pi^*\eta_n=\beta_S(\pi^*[\psi])=0,
\]
which is the assertion under \eqref{eq:UCT}.
\end{proof}


\section{The exact two-filtered row calculation}

The assembly must keep a relation whole while its homogeneous pieces are
being followed.  The bookkeeping device which does this is the relative
Chevalley--Eilenberg row complex
\begin{equation}\label{eq:CE-row-complex}
 \mathcal C_i=\Lambda^i\Rel\otimes U(F),\qquad
 \mathcal C_0=U(F),
\end{equation}
with differential
\begin{align}
 d(\rho_1\wedge\cdots\wedge\rho_i\otimes u)
={}&\sum_{a=1}^i(-1)^{i-a}
 \rho_1\wedge\cdots\widehat\rho_a\cdots\wedge\rho_i
 \otimes\rho_a u \notag\\
 &+\sum_{a<b}(-1)^{a+b}
 [\rho_a,\rho_b]\wedge
 \rho_1\wedge\cdots\widehat\rho_a\cdots
 \widehat\rho_b\cdots\wedge\rho_i\otimes u.
                                                        \label{eq:CE-differential}
\end{align}
It satisfies $d^2=0$, all relation entries in it are full relations, and
\begin{equation}\label{eq:CE-image}
 d(\rho\otimes u)=\rho u,\qquad
 \operatorname{im}(d:\mathcal C_1\to\mathcal C_0)=\Rel U.
\end{equation}

Give a tensor in $\mathcal C_i$ total factor number $i+t$ when its
$U(F)$-entry is a $t$-factor PBW monomial.  The factor-preserving symbol
of \eqref{eq:CE-differential}, in total factor number $p$, is the full
Koszul complex
\begin{equation}\label{eq:CE-Koszul-symbol}
 K_p(\Rel\hookrightarrow F)_i
 =\Lambda^i\Rel\otimes\Sym^{p-i}F.                    
\end{equation}
The lower-factor components are generated, integrally, by
\begin{equation}\label{eq:CE-transfers}
 \begin{split}
 [\cdots|v|u|\cdots]
  &=[\cdots|u|v|\cdots]+[\cdots|[v,u]|\cdots],\\
 [\cdots|v|\rho|\cdots]
  &=[\cdots|\rho|v|\cdots]+[\cdots|[v,\rho]|\cdots].
 \end{split}
\end{equation}
The correction in the second line again has a full relation entry.

To make the collection algorithm integral, filter $\Rel$ by
$\Rel\cap F_{\geq s}$.  The leading-component map embeds
\[
 (\Rel\cap F_{\geq s})/(\Rel\cap F_{\geq s+1})
 \lhook\joinrel\longrightarrow F_s.
\]
Its source is free abelian.  Choose bases successively and lift them;
this gives a triangular basis of $\Rel$ in which a basis relation of
leading weight $s$ has the form
\begin{equation}\label{eq:triangular-relation}
                         \rho=\rho_s+\rho_{>s}.
\end{equation}
A placed $p$-row is a word with $p-1$ homogeneous ordinary entries and
one distinguished basis relation.  During collection the relation in
the row always remains the full $\rho$.  Its active weight edge is the
identity
\begin{equation}\label{eq:active-weight-edge}
 \pr_{\geq s}\rho=\rho_s+\pr_{\geq s+1}\rho;
\end{equation}
the two terms on the right are recorded as components of the quotient
tower, not inserted into the row module as relations.

Iterating the first transfer in \eqref{eq:CE-transfers} gives, for
$B\in F_{\geq2}$ and $x_i\in F_1$, the literal integral identity
\begin{equation}\label{eq:subset-collection}
 Bx_1\cdots x_r
 =\sum_{I\subseteq\{1,\ldots,r\}}
   x_{I^c}[B;x_I],
\end{equation}
where both products and the teeth of the left-normed bracket retain
their original order.  This follows by induction on $r$ from
$Bx=xB+[B,x]$.  In particular the unique one-factor summand is the
$I=\{1,\ldots,r\}$ term, with coefficient $+1$.  Applying
\eqref{eq:subset-collection} simultaneously to all homogeneous
components in \eqref{eq:active-weight-edge} shows why the proper-subset
terms recombine into the full commutator rows required by
\eqref{eq:CE-transfers}.

For the second direction use the tower of quotient maps, rather than
splitting a tail off as a new relation.  The maps
\[
 \Rel\longrightarrow D_k,\quad \rho\longmapsto\pr_{\leq k}\rho,
 \qquad F\longrightarrow A_k
\]
take the factor-$p$ symbol to $K_p(k)$.  The passage from $k$ to $k-1$
is $f_k$, and
\begin{equation}\label{eq:successive-weight-differential}
 \pr_{\leq k}\rho=\pr_{\leq k-1}\rho+\rho_k
\end{equation}
is an equality between two components of the same full-row
differential.  In particular, $\rho_k$ is never asserted to belong to
$\Rel$.

Formally, put $A_{n+1}=F$, $D_{n+1}=\Rel$, and $A_0=D_0=0$.
For fixed $p$, the kernels of the surjections
\[
 K_p(n+1)\longrightarrow K_p(k)
\]
form a finite filtration of the actual Koszul symbol complex
$K_p(\Rel\hookrightarrow F)$; its successive quotient maps are the
$f_k$.  Thus the weight direction below is a filtration of complexes,
not a decomposition of $\rho$ inside $\Rel$.

We now record the finite two-filtered calculation which replaces the
tail-ledger assertion.

\begin{lemma}[closed-square row lemma]\label{lem:closed-square}
Let $\Theta\in\Rel U$ have zero multi-factor PBW coefficient in every
weight at most $2n$, and suppose that its one-factor part is homogeneous
of weight $n+1$.  Exact collection by \eqref{eq:CE-transfers} produces
chains
\begin{equation}\label{eq:packet-chains}
 \chi_k\in K_{m_k}(k)_1\quad(2\leq k<n),
 \qquad \chi_n\in Z_1K_2(n),
\end{equation}
with the following properties.

First, $\partial\chi_k$ is the complete factor-$m_k$ symbol of
$\Theta$ after projection to $\Sym^{m_k}A_k$.  Let
$Q_n(\chi_n)\in C$ denote the oriented primitive coefficient of the
placed factor-two block.  Define the terminal discrepancy, without any
claim that it is a relation tail, by
\begin{equation}\label{eq:terminal-definition}
 \varepsilon:=\iota_C(\pr_{1,n+1}\Theta)
 -\sum_{k=2}^{n-1}\Phi_k(f_k\chi_k)-Q_n(\chi_n),
\end{equation}
where $\iota_C:F_{n+1}\to C$ is the quotient map.  Then
\begin{equation}\label{eq:terminal-discrepancy}
 \varepsilon=-\sum_{k=2}^{n-1}T_k(\partial\chi_k).
\end{equation}
Consequently $\varepsilon=0$.
\end{lemma}

\begin{proof}
Start with any $x\in\mathcal C_1$ satisfying $dx=\Theta$ and expand its
relation entries in the triangular basis
\eqref{eq:triangular-relation}.  Process
factor number in descending order.  At factor number $p\geq3$, order
the ordinary entries and put the distinguished relation at the far
left; at $p=2$ retain the oriented Smith placements used in
Proposition~\ref{prop:quadratic-block}.  Expand a relation into
homogeneous components only to decide the PBW interchanges, and apply
the second identity in \eqref{eq:CE-transfers} simultaneously to the
full relation.  The component not yet reached is therefore a vertical
edge of the quotient tower, not a new relation row.  Append every
full-relation correction to the $(p-1)$-row list.  Within a fixed
factor level, an ordinary transfer lowers the PBW inversion number, a
relation transfer lowers the distance from its prescribed position,
and \eqref{eq:active-weight-edge} raises the active weight.  Hence the
lexicographic measure
\[
 (p,\ \text{unresolved adjacent pairs},\ n+1-s)
\]
strictly decreases at the step being processed.  On the weight
truncation in the hypothesis it is finite, so the algorithm terminates.
Modulo the next factor level the differential is
\eqref{eq:CE-Koszul-symbol}.  Therefore, at $p=m_k$, replacing the full
relation by its image in $D_k$, projecting the other entries to $A_k$,
and forgetting placement gives a chain $\chi_k$ whose Koszul boundary is
exactly the complete projected $p$-factor symbol of $\Theta$.  This
proves the first assertion with integral coefficients.  When $k=n$,
every monomial in $\Sym^2A_n$ has total bracket weight at most $2n$;
the hypothesis on $\Theta$ therefore gives $\partial\chi_n=0$.

It remains to compute the last exposed edge of this finite collection.
Arrange the calculation as a rectangle: factor number decreases from
top to bottom, while the quotient $D_k\to D_{k-1}$ decreases from right
to left.  A cell belonging to $p=m_k$ has upper edge
$\partial\chi_k$.  If $a=\pr_{\leq k-1}\rho$, its horizontal edge is
obtained from \eqref{eq:successive-weight-differential}; since
$b_k(a)=-\rho_k$ in $U_k$, the edge on
$a\otimes x_1\cdots x_{m_k-1}$ is
\begin{equation}\label{eq:cell-horizontal-edge}
 [\rho_k,x_1,\ldots,x_{m_k-1}]
 =\Phi_k(a;x_1\cdots x_{m_k-1}).
\end{equation}
Its vertical edge is obtained by applying $T_k$ to the upper Koszul
boundary.  Formula \eqref{eq:T-k} shows on the same elementary cell that
this is
\begin{equation}\label{eq:cell-vertical-edge}
 T_k\bigl((a+\rho_k)x_1\cdots x_{m_k-1}\bigr)
 =\Theta_k(\bar\rho_k;\bar x_1\cdots\bar x_{m_k-1}).
\end{equation}
Thus the two exposed edges have the same coefficient and opposite
boundary orientations.  This is precisely the chain identity
\eqref{eq:chain-transgression}.

Summing the oriented boundaries of all cells cancels every internal
edge.  This statement involves no division and no saturation: it is
just repeated use of \eqref{eq:CE-transfers} and
\eqref{eq:successive-weight-differential}.  For completeness, the only
critical pairs are the following.  Disjoint transfers commute; three
ordinary entries give ordinary Jacobi; and an overlap involving a
relation gives
\begin{equation}\label{eq:CE-relation-Jacobi}
 [u,[v,\rho]]-[v,[u,\rho]]=[[u,v],\rho].
\end{equation}
All three relation expressions in \eqref{eq:CE-relation-Jacobi} are full
relations.  These are exactly the homogeneous components of $d^2=0$ in
\eqref{eq:CE-row-complex}.  Hence every internal edge cancels once,
independently of the order of collection.

	There is no additional critical pair between transfer and truncation.
	For homogeneous $v\in F_h$, bilinearity and the grading give the literal
	identity
\begin{equation}\label{eq:transfer-truncation-square}
 \pr_{\leq k+h}[v,\rho]=[v,\pr_{\leq k}\rho].
\end{equation}
Thus taking a horizontal quotient edge before or after an adjacent
transfer gives the same entry.  Equations
\eqref{eq:CE-relation-Jacobi} and
	\eqref{eq:transfer-truncation-square}, together with disjoint
	commutation, exhaust the overlaps because every collection step is an
	adjacent interchange or a truncation step.


\medskip
\noindent\emph{Global saturation claim.}
After one-factor, weight-$(n+1)$ projection to $C$, every external
labelled occurrence of the complete two-filtered collection belongs to
exactly one of four families: the outside factor-one edge, a diagonal
horizontal edge, the oriented factor-two edge, or a diagonal vertical
edge.  Their signed sums are respectively
\[
 \iota_C(\pr_{1,n+1}\Theta),\qquad
 -\Phi_k(f_k\chi_k),\qquad -Q_n(\chi_n),\qquad
 +T_k(\partial\chi_k).
\]
We prove this occurrence-level assertion before using it.

A $k$-marked
row is a placed word in which the marked entry
$\langle\rho,k\rangle$ evaluates as $\pr_{\leq k}\rho$; put
$\langle\rho,n+1\rangle=\rho$ and
$\langle\rho,0\rangle=0$.  Besides the first transfer in
\eqref{eq:CE-transfers}, use
\begin{align}
 [\cdots|v|\langle\rho,k\rangle|\cdots]
 &= [\cdots|\langle\rho,k\rangle|v|\cdots]
   +[\cdots|\langle[v,\rho],k+h\rangle|\cdots],
                                                        \label{eq:marked-transfer}\\
 [\cdots|\langle\rho,k\rangle|\cdots]
 &= [\cdots|\langle\rho,k-1\rangle|\cdots]
   +[\cdots|\rho_k|\cdots],
                                                        \label{eq:marked-truncation}
\end{align}
for homogeneous $v\in F_h$, with an index above $n+1$ read as
$n+1$.  These are literal identities: the first follows from
\eqref{eq:transfer-truncation-square}, and the second is
\eqref{eq:successive-weight-differential}.  In particular the
correction in \eqref{eq:marked-transfer} is marked by the full relation
$[v,\rho]\in\Rel$.

Every occurrence, rather than merely every value, will now be retained.
More precisely, label the initial rows by distinct finite words.  When a
labelled occurrence is replaced, append the number of the output term to
its word; if a homogeneous entry or a relation is expanded in a chosen
basis, append also the corresponding basis index.  The integer
coefficient of a child is, by definition, the coefficient of its parent
times the coefficient in the literal replacement.  Thus a child has one
parent even when its value agrees with the value of another child.
Process every child produced at factor number at least three.  At factor
two retain every placement until its image over $A_n$ has been formed.
The resulting complete signed list is exactly the placed realization of
$\chi_n$, and its two-factor symbol is zero by the hypothesis on
$\Theta$.  Every occurrence which can have a primitive descendant of
weight $n+1$ has both inputs in $A_n$, so all target-relevant
occurrences occur in this list.  Pass the whole list, without omission,
through the integral interchanges of
Proposition~\ref{prop:quadratic-block}, and return every one-factor
correction to the list.  Factor-two occurrences outside this projection
cannot contribute to the target projection and are retained as silent
occurrences.  At factor one continue
\eqref{eq:marked-truncation} down to mark zero.  Its homogeneous outputs
are only values in $F$; they are never asserted to lie in $\Rel$.
Ordinary words are PBW-collected completely.  We retain even a child
whose value or eventual projection is zero until that fact has been
checked.

This labelled procedure is finite.  A transfer correction has smaller
factor number, the other transfer output has fewer unresolved adjacent
pairs, and a truncation lowers the mark.  Let $I$ be the signed list of
initial occurrences and let $F_j$ be the list of outputs of the first
$j$ replacements which have not yet been used as inputs.  If
$\partial_{\rm led}s$ denotes the input occurrence of a replacement
minus its signed output list, induction on $j$ gives
\begin{equation}\label{eq:prefix-ledger}
 I-F_j=\sum_{i=1}^{j}\partial_{\rm led}s_i
\end{equation}
in the free abelian group on labelled occurrences.  At termination,
\begin{equation}\label{eq:complete-ledger}
 I-F_N=\sum_{i=1}^{N}\partial_{\rm led}s_i.             
\end{equation}
In particular, an output is either the uniquely labelled input of a
later replacement or is in the final frontier.  This proves at once
that no residual component occurrence can disappear.

We next identify that frontier, still occurrence by occurrence.  Call a
truncation occurrence \emph{$k$-diagonal} if its row has factor number
$m_k$ and, in the primitive bracket descendant being considered, its
homogeneous output $\rho_k$ is the only input in $M$; necessarily all
the other inputs are in $F_1$.  The two children in
\eqref{eq:marked-truncation} have the same parent label and the same
integer coefficient.  The child containing
$\langle\rho,k-1\rangle$ is precisely the corresponding occurrence in
$f_k\chi_k$.  By \eqref{eq:subset-collection}, the child containing
$\rho_k$ has exactly one one-factor descendant, with coefficient $+1$,
namely
\[
                 [\rho_k,x_1,\ldots,x_{m_k-1}].
\]
Consequently the signed sum of all horizontal descendants of the
$k$-diagonal occurrences is exactly
$\Phi_k(f_k\chi_k)$, with the coefficients already present in
$\chi_k$.

There is also a uniquely determined upper partner for each such
occurrence.  Indeed, before the truncation its factor-$m_k$ symbol is
one of the labelled summands of $\partial\chi_k$.  Applying $T_k$ to
that summand gives \eqref{eq:cell-vertical-edge}; applying $\Phi_k$ to
the lower child gives \eqref{eq:cell-horizontal-edge}.  They have the
same coefficient because they have the same parent.  They have opposite
boundary orientations because
\eqref{eq:marked-truncation}, moved to one side, is
\[
 [\cdots|\langle\rho,k\rangle|\cdots]
 -[\cdots|\langle\rho,k-1\rangle|\cdots]
 -[\cdots|\rho_k|\cdots]=0.
\]
Thus the complete $k$-diagonal boundary, with its orientation, is
\begin{equation}\label{eq:labelled-diagonal-boundary}
             -\Phi_k(f_k\chi_k)+T_k(\partial\chi_k).
\end{equation}
This also proves that every lower-context partner and every upper
partner belongs to the same global ledger; neither is supplied by a
separate existence argument.

For completeness, there cannot be a second ancestry for one of these
partners.  Moving a homogeneous entry past a mark and then truncating,
or truncating first and then moving it, gives the two sides of
\eqref{eq:transfer-truncation-square}, with the same inherited label
prefix.  Two transfers with overlapping support give
\eqref{eq:CE-relation-Jacobi}, and disjoint transfers commute.  Hence
every ambiguity is filled by one of the critical squares already listed,
and its two middle occurrences have opposite signs.  Induction on the
finite processing measure shows that the set of cells just described is
downward closed: whenever it contains an input, it contains each
nonsilent output or records that output on the frontier.  This is the
required global, rather than local, saturation assertion.

Here is that induction explicitly.  After any initial segment of the
processing queue, maintain the following three assertions: every
dequeued occurrence has either been paired with its unique parent edge
or recorded on the frontier; every child not yet paired occurs once in
the queue with its inherited coefficient; and whenever one route around
a critical square has been dequeued, the other route is either already
paired or occurs in the queue with the same coefficient.  They hold for
the initial list.  An ordinary transfer preserves them by the first
identity in \eqref{eq:CE-transfers}; a marked transfer preserves them by
\eqref{eq:marked-transfer} and
\eqref{eq:transfer-truncation-square}; a truncation preserves them
because both displayed children in \eqref{eq:marked-truncation} are
enqueued; and an overlapping pair of transfers preserves them by
\eqref{eq:CE-relation-Jacobi}.  The factor-two interchanges are instances
of the same marked transfer, with their factor-one corrections enqueued.
These are all possible processing steps.  Induction therefore proves
the three assertions at termination, when the queue is empty.  This is
also an occurrence-level proof that the coefficients used on the
diagonal and terminal reads are the coefficients of the original global
ledger.

It remains only to classify the recorded frontier.  A one-factor
descendant is a bracket tree in the homogeneous inputs of its ancestor.
Jacobi expresses it integrally as a sum of combs, and the preceding
critical-pair identities retain the labels and coefficients during this
replacement.  A nonzero metabelian comb has at most one input in
$M=F_{\geq2}$.
\begin{itemize}
\item If that input is the relation component $\rho_k$, all other inputs
are in $F_1$.  A comb descended from a factor-$p$ row has weight
$k+p-1$, so weight $n+1$ forces $p=m_k$.  For $2\leq k<n$ this is the
unique $k$-diagonal family above; $k=n$ belongs to the factor-two block;
$k=n+1$ is on the terminal factor-one edge; and $k=1$ is the full
commutator case below.
\item If the unique $M$-input is ordinary, every component of the
relation of weight at least two brackets trivially with it.  Replacing
the weight-one component by the full relation therefore leaves the comb
unchanged, so its image in $C$ is zero.  Two $M$-inputs give zero.
\item If every input has weight one, weight $n+1$ forces $p=n+1$.
Replacing the weight-one relation component by the full relation adds
only terms of weight above $n+1$, and the comb is a full commutator
relation.  It again maps to zero in $C$.  A descendant with a spectator
has more than one PBW factor and is killed by the one-factor projection;
and $p>n+1$ cannot have total weight $n+1$.
\end{itemize}
These cases are disjoint and exhaustive because every descendant has a
unique ancestry label and because every marked truncation produces both
of its children.  In particular this is a classification of all
residual component \emph{occurrences}, not only of the possible values
of isolated combs.

The factor-two boundary uses exactly the signed occurrence list defining
$\chi_n$.  Indeed, no occurrence is removed before that whole list is
put into the canonical Smith placements.  Each interchange has the
literal form $x\rho=\rho x+[x,\rho]$; its correction receives a child
label and is processed at factor one as above.
Proposition~\ref{prop:quadratic-block} therefore reads the primitive contribution of
this same list, including its coefficients and orientation, as
$-Q_n(\chi_n)$.  Thus terminal factor-two rows and the quadratic
primitive read cannot refer to different sets of occurrences.

Finally orient the outside factor-one edge positively and the factor
direction from larger to smaller factor number.  The input-minus-outputs
convention in \eqref{eq:complete-ledger} then gives the following four,
and only four, nonsilent boundary families:
\begin{center}
\begin{tabular}{c|c}
 external family & contribution in $C$ \\ \hline
 terminal one-factor edge &
 $\iota_C(\pr_{1,n+1}\Theta)$ \\
 diagonal horizontal edges, $2\leq k<n$ &
 $-\Phi_k(f_k\chi_k)$ \\
 oriented factor-two edge & $-Q_n(\chi_n)$ \\
 diagonal vertical edges, $2\leq k<n$ &
 $+T_k(\partial\chi_k)$
\end{tabular}
\end{center}
The first entry is the sum of all terminal one-factor component
occurrences after evaluation; exact PBW collection identifies that sum
with $\pr_{1,n+1}\Theta$.  This is not a declaration that any one of
those components is a relation.  An occurrence on a diagonal filtration
cut and its descendant on the outside PBW edge are different labelled
incidences; the intervening critical square pairs them.  Thus the table
does not count a frontier occurrence twice merely because two incidences
have the same value in $F_{n+1}$.  All internal labelled occurrences
cancel in pairs in \eqref{eq:complete-ledger}, every replacement
evaluates to a literal zero, and the exhaustive audit above accounts for
every external occurrence.  Hence
\begin{equation}\label{eq:evaluated-complete-ledger}
 0=\iota_C(\pr_{1,n+1}\Theta)
   -\sum_{k=2}^{n-1}\Phi_k(f_k\chi_k)-Q_n(\chi_n)
   +\sum_{k=2}^{n-1}T_k(\partial\chi_k).
\end{equation}
By \eqref{eq:terminal-definition}, this is
$\varepsilon+\sum_kT_k(\partial\chi_k)=0$, which proves
\eqref{eq:terminal-discrepancy}.

	For clarity, the terminal case $n=2$ is contained in the same ledger
	but can also be read without the notation.  The target has weight three.
	A primitive term coming from three factors is an integral sum of
	$[\rho_1,x,y]$ with $x,y\in F_1$; since $F$ has class three,
	\[
	 [\rho_1,x,y]=[\rho,x,y]\in\Rel,
	\]
	because every higher component of $\rho$ would have weight at least
	four.  More than three positive-weight factors cannot contribute.
	The complete contribution from two factors is, by definition, the
	oriented quadratic block $Q_2(\chi_2)$.  The factor-three full
	commutator leaves are paired, occurrence by occurrence, with their
	marked correction edges by \eqref{eq:prefix-ledger}; the factor-two
	edge is paired with the quadratic block.  There is no diagonal vertical
	edge because there is no $k$ with $2\leq k<2$.  The terminal frontier is
	therefore exactly
	$\iota_C(\pr_{1,3}\Theta)-Q_2(\chi_2)$, and
	\eqref{eq:prefix-ledger} gives
	\[
	 \iota_C(\pr_{1,3}\Theta)-Q_2(\chi_2)=0.
	\]
	This is \eqref{eq:evaluated-complete-ledger} when both intermediate sums
	are empty, and supplies the separate terminal computation explicitly.

Finally, $T_k$ is supported only on
\[
 F_k\Sym^{m_k-1}(F_1)\subseteq\Sym^{m_k}A_k,
\]
of total bracket weight $k+m_k-1=n+1$.  By the first assertion this is
the weight-$(n+1)$, factor-$m_k$ PBW coefficient of $\Theta$, which is
zero.  Thus every summand in
\eqref{eq:terminal-discrepancy} vanishes and $\varepsilon=0$.
\end{proof}

\begin{lemma}[PBW assembly]\label{lem:PBW-assembly}
Let $\theta\in\Rel U$ satisfy \eqref{eq:governing}.  There are chains
$\chi_k$ as in \eqref{eq:packet-chains} such that, in
$C\subseteq\Q/\Z$,
\begin{equation}\label{eq:PBW-normal-form}
 \frac{\zeta}{q}+\Z
 =\sum_{k=2}^{n-1}\Phi_k(f_k\chi_k)
 -\eta_n\bigl(L_1\Sym^2(\pi_n)[\chi_n]\bigr).
\end{equation}
\end{lemma}

\begin{proof}
Apply Lemma~\ref{lem:closed-square} to $\Theta=\theta$.  By
\eqref{eq:terminal-definition} and the conclusion $\varepsilon=0$ of
that lemma,
\[
 \iota_C(\pr_{1,n+1}\theta)
 =\sum_{k=2}^{n-1}\Phi_k(f_k\chi_k)+Q_n(\chi_n).
\]
The occurrence-level saturation proof in that lemma is exactly what
identifies all three terms in this equality; no new frontier
classification is needed here.  Equation \eqref{eq:primitive-witness}
identifies the left side with $\zeta/q+\Z$.
Proposition~\ref{prop:quadratic-block} identifies the oriented quadratic term with
\[
 Q_n(\chi_n)
 =-\eta_n\bigl(L_1\Sym^2(\pi_n)[\chi_n]\bigr).
\]
Substitution gives \eqref{eq:PBW-normal-form}.
\end{proof}

\begin{proposition}[intermediate packet vanishing]
\label{prop:intermediate-vanishing}
For every $2\leq k<n$,
\[
                         \Phi_k(f_k\chi_k)=0.
\]
\end{proposition}

\begin{proof}
By \eqref{eq:chain-transgression},
$\Phi_k(f_k\chi_k)=T_k(\partial\chi_k)$.  The support calculation in
the last paragraph of Lemma~\ref{lem:closed-square} and
\eqref{eq:governing} make the right-hand side zero.
\end{proof}

\section{Proof of the theorem}

\begin{proof}[Proof of Theorem~\ref{thm:main}]
The case $n=1$ was proved in Proposition~\ref{prop:top-reduction}, so
assume $n\geq2$.
By Proposition~\ref{prop:top-reduction}, let
$\zeta z\in\delta_{2n+1}(L)$ and choose the witness
\eqref{eq:witness}.  Lemma~\ref{lem:PBW-assembly} gives
\eqref{eq:PBW-normal-form}.  For every $k<n$,
Proposition~\ref{prop:intermediate-vanishing} makes the $k$th summand zero.
Proposition~\ref{prop:quadratic-vanishing} makes the last summand zero.
Hence
\[
                         \frac\zeta q+\Z=0,
\]
so $q\mid\zeta$ and $\zeta z=0$.
Proposition~\ref{prop:top-reduction} completes the induction.
\end{proof}

\begin{thebibliography}{9}

\bibitem{BM}
L.~Breen and R.~Mikhailov,
\emph{Derived functors of non-additive functors and homotopy theory},
Algebr. Geom. Topol. \textbf{11} (2011), 327--415;
\href{https://arxiv.org/abs/0910.2817}{arXiv:0910.2817}.

\bibitem{PS}
I.~B.~S.~Passi and T.~Sicking,
\emph{Dimension quotients of metabelian Lie rings},
Internat. J. Algebra Comput. \textbf{27} (2017), 251--258;
\href{https://arxiv.org/abs/1602.04919}{arXiv:1602.04919}.

\end{thebibliography}

\end{document}
